\section{Additional results on two-positive maps, continued}

\begin{lemma}[Right-tensor factorization for $VV^\dagger$]
    \label{lem:right_tensor_factorization_mul_adjoint}
    \lean{ChoiJamiolkowski.rightTensorMatrix_mul_conjTranspose_eq_conjTranspose_mul_self}
    \leanok
    \uses{def:choi_right_tensor_factor}
    For every $V\in\MN{D,k}(\C)$, the right-tensor factor satisfies
    $R_{VV^\dagger}=R_V^\dagger R_V$.
\end{lemma}

\begin{proof}
    \leanok
    This is the entrywise calculation obtained by expanding the definition of
    $R_V$ and summing over the middle physical index.
\end{proof}

\begin{lemma}[Factorized right-tensor Choi sandwich]
    \label{lem:factorized_right_tensor_choi_sandwich}
    \lean{ChoiJamiolkowski.rightTensorMatrix_mul_conjTranspose_choi_sandwich_eq}
    \leanok
    \uses{def:choi_right_tensor_factor,
        lem:right_tensor_factorization_mul_adjoint}
    Let $T:\MN{D,D}(\C)\to\MN{D,D}(\C)$ have Choi matrix $\tau$. For every
    $V\in\MN{D,k}(\C)$,
    \begin{align}
        R_{VV^\dagger}\tau R_{VV^\dagger}^\dagger
        &=R_V^\dagger(R_V\tau R_V^\dagger)R_V.
        \label{eq:schwarz_factorized_choi_sandwich}
    \end{align}
\end{lemma}

\begin{proof}
    \leanok
    Replace $R_{VV^\dagger}$ by $R_V^\dagger R_V$ using
    Lemma~\ref{lem:right_tensor_factorization_mul_adjoint} and reassociate the
    matrix products.
\end{proof}

\begin{theorem}[Factorized right projections in the Choi criterion]
    \label{thm:factorized_right_projection_choi_sandwich}
    \lean{ChoiJamiolkowski.IsNPositiveMap.rightTensor_choiMatrix_sandwich_posSemidef_of_mul_conjTranspose}
    \leanok
    \uses{lem:factorized_right_tensor_choi_sandwich,
        thm:k_positive_right_tensor_choi_sandwiches}
    Assume $D>0$. Let $T:\MN{D,D}(\C)\to\MN{D,D}(\C)$ be $k$-positive,
    with Choi matrix $\tau$. If a right-factor matrix has the form
    $P=VV^\dagger$ with $V\in\MN{D,k}(\C)$, then
    $R_P\tau R_P^\dagger\geq0$.
\end{theorem}

\begin{proof}
    \leanok
    By (\ref{eq:schwarz_factorized_choi_sandwich}),
    $R_{VV^\dagger}\tau R_{VV^\dagger}^\dagger
    =R_V^\dagger(R_V\tau R_V^\dagger)R_V$.
    The factor $R_V\tau R_V^\dagger$ is positive by
    Theorem~\ref{thm:k_positive_right_tensor_choi_sandwiches}, and positive
    semidefiniteness is preserved under compression by $R_V$.
\end{proof}

\begin{lemma}[Rank-$k$ Hermitian projection factorization]
    \label{lem:rank_k_hermitian_projection_factorization}
    \lean{Matrix.exists_mul_conjTranspose_of_isHermitian_idempotent_rank}
    \leanok
    Let $P\in\MN{D,D}(\C)$ be Hermitian, idempotent, and of rank $k$.
    Then there exists $V\in\MN{D,k}(\C)$ such that $P=VV^\dagger$.
\end{lemma}

\begin{proof}
    \leanok
    Regard $P$ as a linear operator on $\C^D$. Hermiticity and idempotence say
    that this operator is the orthogonal projection onto its range. Choose an
    orthonormal basis of the range, indexed by $\{0,\ldots,k-1\}$ using the
    rank hypothesis, and let $V$ be the matrix whose columns are these basis
    vectors. The standard rank-one expansion of an orthogonal projection gives
    $P=VV^\dagger$.
\end{proof}

\begin{theorem}[Rank-$k$ right projections in the Choi criterion]
    \label{thm:rank_k_right_projection_choi_sandwich_forward}
    \lean{ChoiJamiolkowski.IsNPositiveMap.rightTensor_choiMatrix_sandwich_posSemidef_of_isHermitian_idempotent_rank}
    \leanok
    \uses{lem:rank_k_hermitian_projection_factorization,
        thm:factorized_right_projection_choi_sandwich}
    Assume $D>0$. Let $T:\MN{D,D}(\C)\to\MN{D,D}(\C)$ be $k$-positive,
    with Choi matrix $\tau$. If $P\in\MN{D,D}(\C)$ is Hermitian, satisfies
    $P^2=P$, and has rank $k$, then $R_P\tau R_P^\dagger\geq0$.
    This is the forward rank-$k$ projection-compression direction of
    \cite[Chapter~3, Proposition~3.1]{Wolf2012Quantum}.
\end{theorem}

\begin{proof}
    \leanok
    By Lemma~\ref{lem:rank_k_hermitian_projection_factorization}, write
    $P=VV^\dagger$ with $V\in\MN{D,k}(\C)$. The conclusion is then exactly
    Theorem~\ref{thm:factorized_right_projection_choi_sandwich}.
\end{proof}

\begin{theorem}[Projection-compression converse for the Choi criterion]
    \label{thm:rank_k_right_projection_choi_sandwich_converse}
    \lean{ChoiJamiolkowski.isNPositiveMap_of_forall_rankProjection_rightTensor_choiMatrix_sandwich_posSemidef}
    \leanok
    \uses{thm:k_positive_iff_right_choi_compressions,
        lem:square_sandwich_fixed_rectangular_compression,
        lem:rank_k_projection_fixing_rectangular_factor}
    Assume $D>0$ and $k\le D$. Let
    $T:\MN{D,D}(\C)\to\MN{D,D}(\C)$ have Choi matrix $\tau$. Suppose that
    $R_P\tau R_P^\dagger\geq0$ for every Hermitian projection
    $P\in\MN{D,D}(\C)$ of rank $k$. Then $T$ is $k$-positive.
\end{theorem}

\begin{proof}
    \leanok
    By Theorem~\ref{thm:k_positive_iff_right_choi_compressions}, it is enough
    to prove positivity of the right-factor compression by an arbitrary
    $X\in\MN{D,k}(\C)$. Choose a rank-$k$ Hermitian projection $P$ with
    $PX=X$ using
    Lemma~\ref{lem:rank_k_projection_fixing_rectangular_factor}. The assumed
    positivity of $R_P\tau R_P^\dagger$, together with
    Lemma~\ref{lem:square_sandwich_fixed_rectangular_compression}, gives
    positivity of the rectangular compression by $X$.
\end{proof}

\begin{theorem}[Rank-$k$ projection form of the Choi criterion]
    \label{thm:rank_k_right_projection_choi_criterion}
    \lean{ChoiJamiolkowski.isNPositiveMap_iff_forall_rankProjection_rightTensor_choiMatrix_sandwich_posSemidef}
    \leanok
    \uses{thm:rank_k_right_projection_choi_sandwich_forward,
        thm:rank_k_right_projection_choi_sandwich_converse}
    Assume $D>0$ and $k\le D$. A linear map
    $T:\MN{D,D}(\C)\to\MN{D,D}(\C)$ is $k$-positive if and only if, for every
    Hermitian projection $P\in\MN{D,D}(\C)$ of rank $k$,
    $R_P\tau R_P^\dagger\geq0$, where $\tau$ is the Choi matrix of $T$.
    This is the rank-$k$ projection formulation of
    \cite[Chapter~3, Proposition~3.1]{Wolf2012Quantum}.
\end{theorem}

\begin{proof}
    \leanok
    The forward direction is
    Theorem~\ref{thm:rank_k_right_projection_choi_sandwich_forward}. The
    reverse direction is
    Theorem~\ref{thm:rank_k_right_projection_choi_sandwich_converse}.
\end{proof}

\begin{definition}[Trace-pairing adjoint]\label{def:trace_pairing_adjoint}
    \lean{Matrix.traceAdjointMap}
    \leanok
    The trace-pairing adjoint $E^* : \MN{D'} \to \MN{D}$ of a linear map
    $E : \MN{D} \to \MN{D'}$ between matrix algebras of possibly different
    dimensions is the adjoint for the bilinear pairing
    $(A,B)\mapsto \tr(AB)$.
\end{definition}

\begin{theorem}[Trace-pairing adjoint identity]\label{thm:trace_pairing_adjoint_identity}
    \lean{Matrix.trace_traceAdjointMap_mul}
    \leanok
    \uses{def:trace_pairing_adjoint}
    For every linear map $E : \MN{D} \to \MN{D'}$, every $\rho \in \MN{D'}$,
    and every $X \in \MN{D}$, one has
    \begin{align}
        \tr(E^*(\rho)X)
        &=\tr(\rho E(X)).
        \label{eq:schwarz_trace_adjoint_pairing}
    \end{align}
\end{theorem}

\begin{proof}
    \leanok
    Write $X=\sum_{i,j}X_{ij}e_{ij}$ in the standard matrix units and use
    linearity in $X$. For $X=e_{ij}$ one has
    $\tr(E^*(\rho)e_{ij})=(E^*(\rho))_{ji}$, while the definition of $E^*$
    gives $(E^*(\rho))_{ji}=\tr(\rho E(e_{ij}))$.
\end{proof}

\begin{theorem}[Trace-pairing adjoint is involutive]\label{thm:trace_pairing_adjoint_involutive}
    \lean{Matrix.traceAdjointMap_traceAdjointMap}
    \leanok
    \uses{def:trace_pairing_adjoint, thm:trace_nondeg}
    For every linear map $E : \MN{D} \to \MN{D'}$, the trace-pairing adjoint
    of $E^*$ is $E$.
\end{theorem}

\begin{proof}
    \leanok
    Fix $X\in\MN{D}$ and an arbitrary matrix $N\in\MN{D'}$. Applying
    (\ref{eq:schwarz_trace_adjoint_pairing}) first to
    $E^* : \MN{D'} \to \MN{D}$ and then to $E$ gives
    \begin{align}
        \tr((E^*)^*(X)N)
        &=\tr(XE^*(N))
        =\tr(E^*(N)X)
        =\tr(NE(X))
        =\tr(E(X)N).
        \notag
    \end{align}
    Both $(E^*)^*(X)$ and $E(X)$ lie in $\MN{D'}$, so, since this holds for
    every $N\in\MN{D'}$, nondegeneracy of the trace pairing on $\MN{D'}$
    gives $(E^*)^*=E$.
\end{proof}

\begin{theorem}[Self-duality of the positive semidefinite cone]
    \label{thm:psd_trace_pairing_self_dual}
    \lean{Matrix.PosSemidef.of_forall_trace_mul_nonneg}
    \leanok
    Let $A \in \MN{D}$ be Hermitian. If $\tr(AB)\geq 0$ for every
    positive semidefinite matrix $B$, then $A\geq 0$.
\end{theorem}

\begin{proof}
    \leanok
    Test the hypothesis on rank-one positive semidefinite matrices
    $B=xx^\dagger$. The identity $\tr(Axx^\dagger)=x^\dagger A x$ and the
    hypothesis give $x^\dagger A x\geq 0$ for every vector $x$, so the
    Hermitian matrix $A$ is positive semidefinite.
\end{proof}

\begin{theorem}[Trace-pairing adjoint of a positive map]\label{thm:trace_pairing_adjoint_positive}
    \lean{IsPositiveMap.traceAdjointMap}
    \leanok
    \uses{def:trace_pairing_adjoint, thm:psd_trace_pairing_self_dual}
    If $E : \MN{D}\to\MN{D}$ is positive, then its trace-pairing adjoint $E^*$
    is positive.
\end{theorem}

\begin{proof}
    \leanok
    Let $\rho\geq 0$. First $E^*(\rho)$ is Hermitian. Indeed, positivity of
    $E$ gives $E(A^\dagger)=E(A)^\dagger$, and hence, for matrix units
    $e_{ij}$,
    \begin{align}
        \overline{\tr(\rho E(e_{ij}))}
        &=\tr(E(e_{ij})^\dagger\rho)
        =\tr(E(e_{ji})\rho)
        =\tr(\rho E(e_{ji})).
        \notag
    \end{align}
    For every $B\geq 0$, (\ref{eq:schwarz_trace_adjoint_pairing}) gives
    $\tr(E^*(\rho)B)=\tr(\rho E(B))$.
    Since $E(B)\geq 0$, the right-hand side is non-negative by positivity of
    the trace product of two positive semidefinite matrices. Self-duality of
    the positive semidefinite cone therefore gives $E^*(\rho)\geq 0$.
\end{proof}

\begin{theorem}[Density representation of positive matrix functionals]
    \label{thm:positive_functional_density_representation}
    \lean{Matrix.exists_density_of_positive_functional}
    \leanok
    Let $f : \MN{D} \to \C$ be a complex-linear functional, where $D>0$. If
    $f(X) \geq 0$ for every $X \geq 0$ and $f(\Id)=1$, then there is a density
    matrix $\rho \in \MN{D}$ such that
    $f(X) = \tr(\rho X)$ for every $X \in \MN{D}$.
\end{theorem}

\begin{proof}
    \leanok
    \uses{thm:trace_pairing_adjoint_identity, thm:trace_pairing_adjoint_positive}
    Define $F : \MN{D} \to \MN{D}$ by $F(X) = f(X)\Id$. For $X \geq 0$ one has
    $f(X) \geq 0$, so $F(X) = f(X)\Id \geq 0$ and $F$ is positive.
    Set $\sigma = D^{-1}\Id$ and $\rho = F^*(\sigma)$. Positivity of the
    trace-pairing adjoint gives $\rho \geq 0$, while
    (\ref{eq:schwarz_trace_adjoint_pairing}) gives
    \begin{align}
        \tr(\rho X)
        &=\tr(\sigma F(X))
        =f(X).
        \notag
    \end{align}
    Taking $X = \Id$ yields $\tr(\rho) = 1$.
\end{proof}

\begin{lemma}[Rank-one domination]
    \label{lem:posSemidef_rankone_domination}
    \lean{Matrix.PosSemidef.eq_nonneg_smul_vecMulVec_of_le_smul_vecMulVec}
    \leanok
    Let $A\in\MN{D}$ be positive semidefinite, let $c\in\C$, and let
    $\psi\in\C^D$. If $A\leq c\ketbra{\psi}{\psi}$, then there is a
    non-negative scalar $a$ such that $A=a\ketbra{\psi}{\psi}$.
\end{lemma}

\begin{proof}
    \leanok
    The assertion is immediate when $\psi=0$. Otherwise set
    $s=\langle\psi,\psi\rangle$, $P=\ketbra{\psi}{\psi}$, and $Q=s^{-1}P$.
    Then $Q$ is the orthogonal projection onto $\C\psi$. Write
    $B=cP-A\geq0$. Since $A+B=cP$ vanishes on the kernel of $Q$, positivity
    implies $A(\Id-Q)=0$.
    Taking adjoints gives $(\Id-Q)A=0$, and therefore $A=QAQ$. Finally,
    \begin{align}
        PAP
        &=\langle\psi,A\psi\rangle P, \notag\\
        A
        &=s^{-2}\langle\psi,A\psi\rangle P.
        \notag
    \end{align}
    The coefficient is non-negative because $A\geq0$ and $s>0$.
\end{proof}

% Scope restriction documented in
% docs/paper-gaps/wolf_prop1_5_one_factor_scope.tex.
\begin{theorem}[Positive retractions onto one matrix factor]
    \label{thm:positive_retraction_one_matrix_factor}
    \lean{Matrix.exists_density_of_positive_retraction_onto_right_factor}
    \leanok
    Let $m,d>0$, and let
    $E:\MN{md}\to\MN{md}$ be a positive complex-linear map. Suppose that the
    image of $E$ is contained in $\Id_m\otimes\MN{d}$ and that
    $E(\Id_m\otimes X)=\Id_m\otimes X$ for every $X\in\MN{d}$. Then there is
    a density matrix $\rho\in\MN{m}$ such that, for every $A\in\MN{md}$,
    \begin{align}
        E(A)
        &=\Id_m\otimes
        \tr_m\!\left((\rho\otimes\Id_d)A\right).
        \label{eq:schwarz_positive_retraction_factor}
    \end{align}
    This is the one-factor case of
    \cite[Proposition~1.5 and Equation~(1.40)]{Wolf2012Quantum}.
\end{theorem}

\begin{proof}
    \leanok
    \uses{lem:posSemidef_rankone_domination,
      thm:positive_functional_density_representation,
      thm:rankone_ray_scalar_rigidity}
    Let $R(A)$ be the right tensor factor of $E(A)$, so that
    $E(A)=\Id_m\otimes R(A)$. For $Y\geq0$ and
    $P_\psi=\ketbra{\psi}{\psi}$, the inequality
    $Y\leq\tr(Y)\Id_m$ and positivity give
    \begin{align}
        0
        &\leq R(Y\otimes P_\psi)
        \leq\tr(Y)P_\psi.
        \notag
    \end{align}
    Hence, for each $v\in\C^d$, there is a scalar $c_v$ such that
    $R(Y\otimes P_v)=c_vP_v$.
    By Theorem~\ref{thm:rankone_ray_scalar_rigidity}, applied to the map
    $X\mapsto R(Y\otimes X)$, there is a scalar $c_Y$ such that
    $R(Y\otimes X)=c_YX$ for every $X\in\MN{d}$.

    Fix a coordinate projection $P_0$ and define
    $f(Y)=R(Y\otimes P_0)_{00}$. Rank-one polarization gives
    $R(Y\otimes X)=f(Y)X$ for all $Y\in\MN{m}$ and $X\in\MN{d}$.
    Positivity of $R$ gives $f(Y)\geq0$ when $Y\geq0$, and the retraction
    property gives $f(\Id_m)=1$. The density representation of positive
    matrix functionals therefore supplies a density matrix $\rho$ satisfying
    $f(Y)=\tr(\rho Y)$.

    Finally, on simple tensors,
    \begin{align}
        \tr_m\!\left((\rho\otimes\Id_d)(Y\otimes X)\right)
        &=\tr(\rho Y)X
        =R(Y\otimes X).
        \notag
    \end{align}
    Since matrix units are simple tensors, linearity proves
    (\ref{eq:schwarz_positive_retraction_factor}) for every
    $A\in\MN{md}$.
\end{proof}

% The remaining fixed-point-projection application is documented in
% docs/paper-gaps/wolf_theorem6_14_fixed_point_projection_gap.tex.
\begin{theorem}[Positive retractions onto finite-dimensional star-algebras]
    \label{thm:positive_retraction_finite_star_algebra}
    \lean{StarSubalgebra.exists_block_densities_of_positive_retraction}
    \leanok
    Let $S\subseteq\MN{D}$ be a unital $*$-subalgebra, and let
    $E:\MN{D}\to\MN{D}$ be a positive complex-linear map whose image lies in
    $S$ and whose restriction to $S$ is the identity. There are positive
    integers $d_k,m_k$, a unitary $U$, and density matrices
    $\rho_k\in\MN{m_k}$ such that
    \begin{align}
        U^*SU
        &=\bigoplus_k \Id_{m_k}\otimes\MN{d_k},
        \label{eq:schwarz_star_algebra_blocks}
    \end{align}
    and, for every $A\in\MN{D}$,
    \begin{align}
        E(A)
        &=U\left(\bigoplus_k \Id_{m_k}\otimes
        \tr_{m_k}\!\left((\rho_k\otimes\Id_{d_k})(U^*AU)_{kk}\right)\right)U^*.
        \label{eq:schwarz_star_algebra_retraction}
    \end{align}
    Equivalently, cyclicity of the partial trace permits the factor
    $\rho_k\otimes\Id_{d_k}$ to be placed on the right of $(U^*AU)_{kk}$.
    This is \cite[Proposition~1.5 and Equation~(1.40)]{Wolf2012Quantum}.
\end{theorem}

\begin{proof}
    \leanok
    \uses{thm:positive_retraction_one_matrix_factor,
      thm:starSubalgebra_unitary_block_diagonal_iff}
    Conjugate by $U$ and use (\ref{eq:schwarz_star_algebra_blocks}) to
    identify the ambient space with
    $\bigoplus_k\C^{m_k}\otimes\C^{d_k}$. Let $P_k$ be the central projection
    onto the $k$-th summand and let $T_k$ be the $k$-th diagonal block of the
    conjugated map. Since $T_k(\Id-P_k)=0$, positivity implies
    $T_k((\Id-P_k)A)=T_k(A(\Id-P_k))=0$.
    Hence $T_k(A)=T_k(P_kAP_k)$. The restriction to the $k$-th diagonal
    summand is a positive retraction onto
    $\Id_{m_k}\otimes\MN{d_k}$, so the one-factor theorem gives the density
    matrix $\rho_k$. Taking the direct sum and conjugating back proves
    (\ref{eq:schwarz_star_algebra_retraction}). Entrywise expansion of the
    partial trace gives
    \begin{align}
        \tr_{m_k}((\rho_k\otimes\Id)A_{kk})
        &=\tr_{m_k}(A_{kk}(\rho_k\otimes\Id)).
        \notag
    \end{align}
\end{proof}

\begin{theorem}[Trace-pairing adjoints preserve positivity exactly]\label{thm:trace_pairing_adjoint_positive_iff}
    \lean{isPositiveMap_traceAdjointMap_iff}
    \leanok
    \uses{thm:trace_pairing_adjoint_positive, thm:trace_pairing_adjoint_involutive}
    A linear map is positive if and only if its trace-pairing adjoint is
    positive.
\end{theorem}

\begin{proof}
    \leanok
    One direction is Theorem~\ref{thm:trace_pairing_adjoint_positive}. The
    other follows by applying the same theorem to $E^*$ and using
    Theorem~\ref{thm:trace_pairing_adjoint_involutive}.
\end{proof}

\begin{theorem}[Trace-pairing adjoint of an ampliation]
    \label{thm:trace_pairing_adjoint_ampliation}
    \lean{nPositiveAmpliation_traceAdjointMap}
    \leanok
    \uses{def:blockwise_ampliation, def:trace_pairing_adjoint,
        thm:trace_pairing_adjoint_identity}
    For every linear map $E : \MN{D}\to\MN{D}$, the trace-pairing adjoint
    commutes with ampliation:
    \begin{align}
        (E^{(k)})^*
        &=(E^*)^{(k)}.
        \label{eq:schwarz_trace_adjoint_ampliation}
    \end{align}
\end{theorem}

\begin{proof}
    \leanok
    For block matrices write $\rho_{pq}$ and $X_{pq}$ for their
    $\MN{D}$-valued blocks. Expanding the trace by blocks and applying
    (\ref{eq:schwarz_trace_adjoint_pairing}) gives
    \begin{align}
        \tr((E^*)^{(k)}(\rho)X)
        &=\sum_{p,q}\tr(E^*(\rho_{pq})X_{qp})
        =\sum_{p,q}\tr(\rho_{pq}E(X_{qp}))
        =\tr(\rho E^{(k)}(X)).
        \notag
    \end{align}
    Nondegeneracy of the trace pairing identifies the two adjoints, proving
    (\ref{eq:schwarz_trace_adjoint_ampliation}).
\end{proof}

\begin{theorem}[Trace-pairing adjoint of a $k$-positive map]
    \label{thm:k_positive_trace_pairing_adjoint}
    \lean{IsNPositiveMap.traceAdjointMap}
    \leanok
    \uses{def:k_positive_map, thm:trace_pairing_adjoint_ampliation,
        thm:trace_pairing_adjoint_positive}
    If $E$ is $k$-positive, then $E^*$ is $k$-positive.
\end{theorem}

\begin{proof}
    \leanok
    Since $E$ is $k$-positive, its ampliation $E^{(k)}$ is positive. The
    trace-pairing adjoint of this positive map is positive, and
    (\ref{eq:schwarz_trace_adjoint_ampliation}) identifies that adjoint with
    $(E^*)^{(k)}$.
\end{proof}

\begin{theorem}[$k$-positivity is trace-adjoint invariant]
    \label{thm:k_positive_trace_pairing_adjoint_iff}
    \lean{isNPositiveMap_traceAdjointMap_iff}
    \leanok
    \uses{thm:k_positive_trace_pairing_adjoint,
        thm:trace_pairing_adjoint_involutive}
    A linear map is $k$-positive if and only if its trace-pairing adjoint is
    $k$-positive.
\end{theorem}

\begin{proof}
    \leanok
    One direction is
    Theorem~\ref{thm:k_positive_trace_pairing_adjoint}. The other follows by
    applying the same theorem to $E^*$ and using
    Theorem~\ref{thm:trace_pairing_adjoint_involutive}.
\end{proof}

\begin{theorem}[Completely positive maps are $k$-positive]
    \lean{IsCPMap.isNPositiveMap}
    \leanok
    \uses{def:cp_map}
    Every completely positive map is $k$-positive for all $k \ge 1$.
\end{theorem}

\begin{theorem}[One-positive maps are positive]
    \label{thm:one_positive_iff_positive}
    \lean{isNPositiveMap_one_iff_isPositiveMap}
    \leanok
    \uses{def:k_positive_map, def:positive_map}
    A linear map is $1$-positive if and only if it is positive.
\end{theorem}

\begin{proof}
    \leanok
    The ampliation by $\MN{1}$ is identified with the original matrix algebra:
    all block indices have the unique value, so positivity of
    $E\otimes\id_1$ is exactly positivity of $E$.
\end{proof}

\begin{theorem}[$D$-positive maps on nonzero $\MN{D}$ are completely positive]
    \label{thm:d_positive_iff_completely_positive}
    \lean{isCPMap_iff_isNPositiveMap_card}
    \leanok
    \uses{def:k_positive_map, def:cp_map, thm:cp_iff_choi_posSemidef}
    For maps $E:\MN{D}\to\MN{D}$ with $D\ge 1$, complete positivity is
    equivalent to $D$-positivity.
\end{theorem}

\begin{proof}
    \leanok
    Complete positivity implies $k$-positivity for every $k$. Conversely,
    assume $D\ge 1$. If $E$ is $D$-positive, apply $E\otimes\id_D$ to the
    maximally entangled projector $\ketbra{\Omega}{\Omega}$. The result is
    the Choi matrix of $E$, hence it is positive semidefinite. Choi's theorem
    then gives complete positivity.
\end{proof}

\begin{theorem}[$(k\!+\!1)$-positive implies $k$-positive]
    \lean{IsNPositiveMap.mono}
    \leanok
    Positivity under amplification is monotone in $k$.
\end{theorem}

\begin{theorem}[$k$-positive implies $m$-positive for $m \le k$]
    \lean{IsNPositiveMap.mono_of_le}
    \leanok
    If a linear map is $k$-positive and $m \le k$, then it is $m$-positive.
\end{theorem}

\begin{theorem}[Zero map is $k$-positive]
    \lean{IsNPositiveMap.zero}
    \leanok
    The zero map is $k$-positive.
\end{theorem}

\begin{theorem}[Sums of $k$-positive maps]
    \lean{IsNPositiveMap.add}
    \leanok
    If $E$ and $F$ are $k$-positive, then $E+F$ is $k$-positive.
\end{theorem}

\begin{theorem}[Nonnegative multiples of $k$-positive maps]
    \lean{IsNPositiveMap.smul_nonneg}
    \leanok
    If $E$ is $k$-positive and $c \ge 0$, then $cE$ is $k$-positive.
\end{theorem}

\begin{theorem}[Closedness of $k$-positive maps]
    \lean{isClosed_setOf_isNPositiveMap}
    \leanok
    For each $k$, the set of $k$-positive linear maps is closed in the
    finite-dimensional space of linear maps.
\end{theorem}

\subsection{Douglas-type factorization lemmas}

\begin{theorem}[Range inclusion implies right factorization]
    \lean{Douglas.factorization_of_range_mulVecLin_le}
    \leanok
    If $\ran(A) \subseteq \ran(B)$, then
    there exists $C$ such that $A = B C$.
\end{theorem}

\begin{proof}
    In finite dimensions, $\ran(A) \subseteq \ran(B)$ implies
    $A=B(B^\dagger B)^\dagger B^\dagger A$, where
    $(B^\dagger B)^\dagger$ denotes the Moore--Penrose pseudoinverse.
    Set $C=(B^\dagger B)^\dagger B^\dagger A$.
\end{proof}

\begin{theorem}[Vectorwise range criterion implies factorization]
    \lean{Douglas.factorization_of_forall_mulVec_mem_range}
    \leanok
    If $A v \in \ran(B)$ for every vector $v$, then
    there exists $C$ with $A = B C$.
\end{theorem}

\begin{theorem}[Hilbert--Schmidt contraction]\label{thm:hs_contraction}
    \lean{KadisonSchwarz.hilbertSchmidt_contraction}
    \leanok
    \uses{def:unital_kraus, def:tp_kraus}
    If a Kraus map is both unital and trace-preserving, then
    $\tr(E(X)^\dagger E(X)) \le \tr(X^\dagger X)$
    for all $X \in \MN{D}$.
\end{theorem}

\begin{proof}\leanok\uses{thm:kadison_schwarz}
    By (\ref{eq:schwarz_kadison}),
    $E(X^\dagger X) - E(X)^\dagger E(X) \ge 0$, so
    $\tr(E(X^\dagger X)) \ge \tr(E(X)^\dagger E(X))$.
    Trace preservation gives $\tr(E(X^\dagger X)) = \tr(X^\dagger X)$.
\end{proof}

\begin{lemma}[Trace pairing for Kraus map and adjoint map]\label{lem:trace_mul_map_eq_trace_adjointMap_mul}
    \lean{Kraus.trace_mul_map_eq_trace_adjointMap_mul}
    \leanok
    \uses{def:kraus_map, def:kraus_adjoint_map}
    For all $X, Y \in \MN{D}$,
    $\tr\!(X\,E(Y)) = \tr\!(E^*(X)\,Y)$.
\end{lemma}

\begin{lemma}[Positive-definite trace test]\label{lem:posSemidef_eq_zero_of_posDef_trace_mul_eq_zero}
    \lean{Kraus.posSemidef_eq_zero_of_posDef_trace_mul_eq_zero}
    \leanok
    If $A > 0$, $X \ge 0$, and $\tr(A X) = 0$, then $X = 0$.
\end{lemma}

\begin{theorem}[Fixed-point peripheral eigenvectors satisfy KS equality]
    \label{thm:ks_equality_of_peripheral_eigenvector_of_fixedPoint}
    \lean{Kraus.ks_equality_of_peripheral_eigenvector_of_fixedPoint}
    \leanok
    \uses{thm:kadison_schwarz, lem:trace_mul_map_eq_trace_adjointMap_mul,
        lem:posSemidef_eq_zero_of_posDef_trace_mul_eq_zero}
    Let $E$ be unital and let $E^*$ be trace-preserving with a positive
    definite fixed point $\rho > 0$. If $E(X) = \mu X$ with $|\mu| = 1$, then
    $E(X^\dagger X) = E(X)^\dagger E(X)$.
\end{theorem}
